\documentclass[10pt]{article}
% Math Packages
\usepackage{amsmath, mathtools}
\usepackage{amssymb}
\usepackage{amsthm}
\usepackage{amsfonts}
\usepackage{bbm}
\usepackage{breqn}
\usepackage[margin=1in]{geometry}
\usepackage{graphicx}
\usepackage{tikz}
\usetikzlibrary{arrows.meta}
\usetikzlibrary{calc}
\usepackage{forest}
\usepackage{tikz-qtree}
\graphicspath{ {./images/} }
\usepackage{hyperref}
\usepackage[capitalize]{cleveref}
\usepackage[shortlabels]{enumitem}
\usetikzlibrary{arrows,matrix,positioning}
\usepackage{multicol}

% for the pipe symbol
\usepackage[T1]{fontenc}

% Citing theorems by name. (source: https://tex.stackexchange.com/questions/109843/cleveref-and-named-theorems)
\makeatletter
\newcommand{\ncref}[1]{\cref{#1}\mynameref{#1}{\csname r@#1\endcsname}}

\def\mynameref#1#2{%
  \begingroup
  \edef\@mytxt{#2}%
  \edef\@mytst{\expandafter\@thirdoffive\@mytxt}%
  \ifx\@mytst\empty\else
  \space(\nameref{#1})\fi
  \endgroup
}
\makeatother

% Colorful Notes
\usepackage{color} \definecolor{Red}{rgb}{1,0,0} \definecolor{Blue}{rgb}{0,0,1}
\definecolor{Purple}{rgb}{.5,0,.5} \def\red{\color{Red}} \def\blue{\color{Blue}}
\def\gray{\color{gray}} \def\purple{\color{Purple}}
\newcommand{\rnote}[1]{{\red [#1]}} % \rnote{foo} gives '[foo]' in red
\newcommand{\pnote}[1]{{\purple [#1]}} % \pnote{foo} gives '[foo]' in purple
\newcommand{\bnote}[1]{{\blue #1}} % \bnote{foo} gives 'foo' in blue
\newcommand{\gnote}[1]{{\gray #1}} % \gnote{foo} gives 'foo' in gray
\newcommand{\Max}[1]{{\purple [#1]}} % \bnote{foo} then 'foo' is blue


% Claim numbering (the counter restarts after each proof environment)
\newcounter{claimcount}
\setcounter{claimcount}{0}
\newenvironment{claim}{\refstepcounter{claimcount}\par\addvspace{\medskipamount}\noindent\textbf{Claim \arabic{claimcount}:}}{}
\usepackage{etoolbox}
\AtBeginEnvironment{proof}{\setcounter{claimcount}{0}}
\newenvironment{claimproof}{\par\addvspace{\medskipamount}\noindent\textit{Proof of Claim  \arabic{claimcount}.}}{\hfill\ensuremath{\qedsymbol} \tiny{Claim}

  \medskip}
% Add claim support to cleverref
\crefname{claimcount}{Claim}{Claims}


% Math Environments
\newtheorem{theorem}{Theorem}
\newtheorem{assumption}[theorem]{Assumption}
\newtheorem{lemma}[theorem]{Lemma}
\newtheorem{proposition}[theorem]{Proposition}
\newtheorem{corollary}[theorem]{Corollary}
\newtheorem{question}[theorem]{Question}
\theoremstyle{definition}
\newtheorem{definition}[theorem]{Definition}
\newtheorem{remark}[theorem]{Remark}
\newtheorem{example}[theorem]{Example}
\newtheorem{notation}[theorem]{Notation}
\newtheorem{problem}[theorem]{Problem}

% Matrices and Column Vectors. 
\usepackage{stackengine}
\setstackgap{L}{1.0\normalbaselineskip}
\usepackage{tabstackengine}
\setstackEOL{;}% row separator
\setstackTAB{,}% column separator
\setstacktabbedgap{1ex}% inter-column gap
\setstackgap{L}{1.5\normalbaselineskip}% inter-row baselineskip
\let\nmatrix\bracketMatrixstack  %Usage: \nmatrix{1,2,3\4,5,6}
\newcommand\cv[1]{\setstackEOL{,}\bracketMatrixstack{#1}} %usage: \cv{1,2,3}

% table alignment
\usepackage{array}
\usepackage{booktabs}

% Custom Math Coqmmands
\newcommand{\vt}{\vskip 5mm} % vertical space
\newcommand{\fl}{\noindent\textbf} % first line
\newcommand{\Fl}{\vt\noindent\textbf} % first line with space above
\newcommand{\norm}[1]{\left\lVert#1\right\rVert} % norm
\newcommand{\pnorm}[1]{\left\lVert#1\right\rVert_p} % p-norm
\newcommand{\qnorm}[1]{\left\lVert#1\right\rVert_q} % q-norm
\newcommand{\1}[1]{\textbf{1}_{\left[#1\right]}} % indicator function 
\def\limn{\lim_{n\to\infty}} % shortcut for lim as n-> infinity
\def\sumn{\sum_{n=1}^{\infty}} % shortcut for sum from n=1 to infinity
\def\sumkn{\sum_{k=1}^{n}} % shortcut for sum from k=1 to n
\def\sumin{\sum_{i=1}^{n}} % shortcut for sum from i=1 to n
\def\SAs{\sigma\text{-algebras}} % shortcut for $\sigma$-algebras
\def\SA{\sigma\text{-algebra}} % shortcut for $\sigma$-algebra
\def\Ft{\mathcal{F}_t} % time-indexed sigma-algebra (t)
\def\Fs{\mathcal{F}_s} % time-indexed sigma-algebra (s)
\def\F{\mathcal{F}} % sigma-algebra
\def\G{\mathcal{G}} % sigma-algebra
\def\R{\mathbb{R}} % Real numbers
\def\N{\mathbb{N}} % Natural numbers
\def\Z{\mathbb{Z}} % Integers
\def\E{\mathbb{E}} % Expectation
\def\P{\mathbb{P}} % Probability
\def\Q{\mathbb{Q}} % Q probability
\def\dist{\text{dist}} %Text 'dist' for things like 'dist(x,y)'
\newcommand{\indep}{\perp \!\!\! \perp}  %independence symbol
\def\Var{\mathrm{Var}} % Variance
\def\tr{\mathrm{tr}} % trace

% Brackets and Parentheses
\def\[{\left [}
    \def\]{\right ]}
% \def\({\left (}
%   \def\){\right )}



\usepackage{color}
\definecolor{Red}{rgb}{1,0,0}
\definecolor{Blue}{rgb}{0,0,1}
\definecolor{Purple}{rgb}{.5,0,.5}
\def\red{\color{Red}}
\def\blue{\color{Blue}}
\def\gray{\color{gray}}
\def\purple{\color{Purple}}
\definecolor{RoyalBlue}{cmyk}{1, 0.50, 0, 0}
\newcommand{\dempfcolor}[1]{{\color{RoyalBlue}#1}} 
\newcommand{\demph}[1]{\dempfcolor{\textit{\textbf{#1}}}}

% comment exactly one of the following line to show / hide the solutions
% \newcommand{\solution}[1]{{\purple #1}} % uncomment to show the solutions
\newcommand{\solution}[1]{} % uncomment to hide the solutions

\usepackage{pst-poker}


\begin{document}
\begin{center}
  \section*{Math 372: Worksheet 3}
  \subsection*{February 20, 2026}
\end{center}
\textit{\textbf{Instructions:} Write your answers on a separate sheet.}

\begin{problem}
  Suppose we lay out the following four playing cards in a line like this:
  \begin{center}
    \crdAs \crdtwos \crdtres \crdfours
  \end{center}
  A \demph{random transposition} is where two different cards are selected at random and
  their positions swapped. For example, the transposition $(23)$ swaps cards 2
  and 3. If the cards started in the above order, then the effect of applying
  $(23)$ would be to reorder the cards as follows:
  \begin{center}
    \crdAs \crdtres \crdtwos \crdfours
  \end{center}

  
  Anna's nemesis Mark challenges her to a betting game. In the game, Mark lays
  out four cards in the order $1,2,3,4$ and then performs $k$ random
  transpositions, one after the other. As this happens, Anna is blindfolded,
  so she can't see which transpositions were performed---however, because she
  can hear Mark moving the cards, she does know the value of $k$. Anna then
  has to guess the final order of the cards. If she guesses correctly, Mark
  will pay her $\$20$; otherwise she has to pay Mark $\$1$.

  Mark believes that as long as he performs at least $4$ transpositions, then
  all $4!=24$ permutations of the cards will be equally likely, so that Anna
  is expected to lose money when she plays the game. Moreover Mark believes
  that Anna is a sucker and he can make lots of money from her.

  But Mark is wrong. Design a guessing strategy for Anna to bleed Mark for all
  he's worth.

  \fl{Hint:} \textit{First consider the simpler problem where the game is
    played with just 2 cards, and design a guessing strategy for Anna in that
    case.}

  \textit{Then, consider the game with 3 cards. In that case, there are six
    outcomes (shown below). For each $k=2,3,4,5$, write down a table which
    gives the probability mass function these six outcomes. This will lead you
    to to a good guessing strategy for Anna.}
  \begin{align*}
    &\scalebox{0.4}{\crdAs \crdtwos \crdtres}  &\scalebox{0.4}{\crdAs \crdtres \crdtwos} \\
    &\scalebox{0.4}{\crdtwos \crdAs \crdtres}  &\scalebox{0.4}{\crdtres \crdtwos \crdAs} \\
    &\scalebox{0.4}{\crdtwos \crdtres \crdAs}  &\scalebox{0.4}{\crdtres \crdAs \crdtwos} 
  \end{align*}
  \textit{Once you understand the case with 3 cards, move to
    4.}
\end{problem}

\begin{problem}
  What does Mark think his probability of winning is? What does he think his
  expected payoff for playing the game is?

  \textit{Hint:} use the formula
  \begin{equation*}
    \text{expected payoff} = \text{(payoff if he wins)}\P\left[\text{win} \right] +
    (\text{payoff if he loses}) \P\left[\text{lose} \right]
  \end{equation*}
  
\end{problem}

\begin{problem}
  What is Anna's expected payoff if she adopts your strategy? If she adopts
  your strategy, how much money can she expect to make if she agrees to play
  the game 100 times?
\end{problem}

\begin{problem}[Exercise 3.2.16]
  For the last 21 months, the Ukrainian city of Kherson has been subject to an
  ongoing campaign of drone attacks by Russian forces entrenched across the
  nearby Dnipro River. Suppose that $25\%$ of Ukrainian air-defense crews in
  the city are equipped with an experimental electromagnetic drone jammer.

  Suppose that, on a given night, four crews are selected at random to defend
  a neighborhood near the riverbank. Let $X$ denote the number among the four
  that are equipped with the jammer.

  \begin{enumerate}[(a)]
    \item Find the probability distribution of $X$. \textit{[Hint: Let $J$
    denote a crew equipped with the jammer and $N$ one who does not. Then one
    possible outcome is \texttt{JJNJ}, with probability $(.25)(.75)(.25)(.25)$
    and associated $X$ value $3$. There are 15 other outcomes.]}
    \item Draw the corresponding probability histogram.
    \item What is the most likely value for $X$?
    \item What is the probability that at least two of the four selected crews
    have the jammer?
  \end{enumerate}
\end{problem}

\begin{problem}[Excersise 3.2.21 in the textbook]
  Suppose that you read through this year’s issues of the New York Times and
  record each number that appears in a news article---the income of a CEO, the
  number of cases of wine produced by a winery, the age of a
  celebrity, and so on. Now focus on the leading digit of each number, which
  could be $1, 2, \ldots , 8$, or $9$. Your first thought might be that the
  leading digit $X$ of a randomly selected number would be equally likely to
  be one of the nine possibilities. However, much empirical evidence as well
  as some theoretical arguments suggest an alternative probability
  distribution called Benford's law:
  \begin{equation*}
    p(x) = \P\left[\text{first digit is }x \right] = \log_{10}\left(
      \frac{x+1}{x} \right), \quad \text{for }x=1,2,\ldots,9.
  \end{equation*}
  \begin{enumerate}[(a)]
    \item Without computing individual probabilities from this formula, show
    that it specifies a legitimate pdf (i.e., show that $\sum_{k=1}^{9}p(k)=1$
    and that $p(x) \geq 0$ for all $x=1,\ldots,9$).
    \item Now compute the individual probabilities and compare to the
    corresponding discrete distribution.
    \item Plot the cdf of $X$.
    \item Using the cdf, what is the probability that the leading digit is at
    most 3? At least 5?
  \end{enumerate}
  \textit{[Note: Benford's law is the basis for some auditing procedures used
  to detect fraud in financial reporting---for example, by the Internal
  Revenue Service]}
\end{problem}


For the remaining problem, we'll need the following definition.
\begin{definition}
  Let $X$ be a discrete random variable with finite expected value
  $\mu= \E\left[X \right]$. The \demph{variance} of $X$, denoted $\Var(X)$, is
  the quantity
  \begin{equation*}
    \Var(X) = \sum_{x} \left( x- \mu\right)^{2}  
  \end{equation*}
  where the sum ranges over all distinct values $x$ that the random variable
  $X$ can take.
\end{definition}

\begin{problem}
  A certain brand of upright freezer is available in three different rated
  capacities: $16$,  $18$, and $20$ cubic feet. Let $X$ be the rated capacity
  of a freezer of this brand sold at a certain store. Suppose that $X$ has
  probability mass function
  \begin{center}
    \begin{tabular}[c]{c|ccc}
      $x$&16&18&20\\ 
      \hline
      $p(x)$&.2&.5&.3 
    \end{tabular}
  \end{center}
  \begin{enumerate}[(a)]
    \item Compute $\E\left[X \right], \E\left[X^{2} \right], \Var(X)$.
    \item If the price of a freezer having capacity $X$ is $70X-650$, what is
    the expected price paid by the next customer to buy a freezer?
    \item What is the variance of the price paid by the next customer?
    \item Suppose that although the rated capacity of a freezer is $X$, the
    actual capacity is $h(X)=X-.008X^{2}$. What is the expected actual
    capacity of the freezer purchased by the next customer.
  \end{enumerate}
\end{problem}


% \begin{problem}
%   Compute $\E\left[X \right]$, $\E\left[X^{2} \right]$ and $\Var(X)$ for the
%   following random variables:
%   \begin{enumerate}[(a)]
%     \item $X = (\text{value of one roll of a 6-sided dice})$.
%     \item A fair toss. $Y = \left\{ \begin{array}{l@{\quad:\quad}l} 1 &  \text{with
%           probability }\frac{1}{2}\\0 & \text{with probability }\frac{1}{2} \end{array}\right.$
%     \item A biased coin toss. $Z = \left\{ \begin{array}{l@{\quad:\quad}l} 1 &  \text{with
%           probability }p\\0 & \text{with probability }1-p \end{array}\right.$
%   \end{enumerate}
% \end{problem}

\end{document}